\section{Introduction}
%Try to follow:

%- molecular optimization via discrete diffusion

%- unconditional gen. is not enough. we care about task-relevant designs hence use test-time training, where we can adapt the generative model via online feedback

%- problem: online feedback is very expensive

%- research question: how to optimally query and use costly online feedback for TT training?

%- We consider several important design choices of this process (exploration, objective missmatch etc. see slack -- present clearly \textbf{what they do intuitively}) and analyse them rigorously.

%- Results..

Discrete diffusion models have become a powerful tool for molecular generation, such as for small molecules \citep{leeGenMolDrugDiscovery2025, schiffSimpleGuidanceMechanisms2025} and proteins \citep{wangDiffusionLanguageModels2024}, because they can model realistic chemical structures directly while producing diverse candidate molecules. 
In molecular design, however, unconditional generation is rarely the final goal. 
A pretrained generator encodes a broad molecular prior whose probability mass is biased toward the pretraining distribution, while the design task defines a different target: molecules that score highly under a task-specific oracle. The generator must therefore be steered toward task-relevant regions of chemical space without losing the molecular validity that made the pretrained prior useful in the first place \citep{HowArtificialIntelligence,  ueharaInferenceTimeAlignmentDiffusion2025}.
A common way to do so is test-time adaptation, where one iteratively samples molecules, evaluates a small subset with a task oracle, and uses the feedback to steer future generations \citep{yangSteeringGenerativeModels2025, stoccoGuidingGenerativeModels2025}.
%A pretrained generator must produce parseable SMILES corresponding to chemically valid molecues, and beyond that score well for a specific task: binding to a target, improving a protein fitness score, satisfying a docking pipeline, or eventually passing a wet-lab assay. A common way to specialize a generator to such a task is test-time adaptation: sample molecules, evaluate them with a task oracle, and use the feedback to steer future generations.
%A common way to do so is test-time adaptation: sample molecules, evaluate a small subset with a task oracle, and use the feedback to steer future generations.





%Molecular optimization using generative models \citep{HowArtificialIntelligence} is often bottlenecked by feedback. A generative model can propose many candidate molecules, but each candidate becomes useful only after it is scored by an oracle, such as a learned binding-affinity model \citep{jiangFlashAffinityBridgingAccuracySpeed}, a structure-prediction/docking pipeline \citep{passaro2025boltz2}, or ultimately a wet-lab assay. These oracles are often expensive, non-differentiable and noisy relative to sampling from the generator. As a result, the central test-time question is not only how good the pretrained generator is, but how to allocate a limited number of oracle calls and subsequent computation \citep{didi2026scaling}.


%One attractive but relatively underexplored direction is test-time training, or online fine-tuning, for generative molecule optimization \citep{rector-brooksSteeringMaskedDiscrete2024, suIterativeDistillationRewardGuided2025, ueharaRewardGuidedIterativeRefinement2025, yangSteeringGenerativeModels2025, stoccoGuidingGenerativeProtein2025, stoccoGuidingGenerativeModels2025}. Test-time training uses evaluated molecules to update the generator itself. This has a natural advantage when oracle calls are scarce: each evaluation can influence the distribution of future samples, rather than only the sample selected now. In this paper, we focus on online fine-tuning of discrete diffusion models and ask which additions make this feedback loop work best.

%The difficulty is that task feedback is often the bottleneck. A generator can sample many molecules cheaply, but it cannot determine which samples are task-relevant until some of them are evaluated by an oracle, such as a learned binding-affinity model \citep{jiangFlashAffinityBridgingAccuracySpeed}, a structure-prediction or docking pipeline \citep{passaro2025boltz2}, or a wet-lab assay. These evaluations can be expensive relative to sampling from the generator. As a result, the central test-time question is not only how strong the pretrained generator is, but how to allocate a limited number of oracle calls and subsequent computation \citep{didi2026scaling}. In this paper, we ask how the full online training loop should use costly oracle feedback when optimizing a discrete diffusion model at test time.

The difficulty with  test-time adaptation is that it leads to a coupled design problem. A generator can sample many molecules cheaply, but it cannot determine which samples are task-relevant until some are evaluated by an oracle, such as a learned binding-affinity model \citep{jiangFlashAffinityBridgingAccuracySpeed}, a structure-prediction or docking pipeline \citep{passaro2025boltz2}, or a wet-lab assay. These evaluations can be expensive relative to sampling from the generator, so each oracle call must serve two roles: selecting useful candidate molecules now and providing feedback that can improve future generations. A key question is therefore not only how strong the pretrained generator is, or which fine-tuning objective is used, but how to reason holistically about the full online training loop: how it should allocate, reuse, and transform a limited stream of oracle feedback \citep{didi2026scaling}.



%We study this question through online fine-tuning for discrete-diffusion molecular optimization. Starting from a pretrained generator, the optimizer repeatedly samples molecules, decides which candidates to evaluate, receives oracle scores, and updates the generator from the scored molecules. This loop creates a growing history of feedback, but it leaves two central design problems open. The first is exploration: the optimizer must both query informative molecules to learn the unknown reward landscape and generate diverse candidates beyond the model's current high-probability regions. The second is objective mismatch: the oracle rewards only task-relevant molecules, while likelihood-based fine-tuning can spend capacity modeling low-value or invalid samples. These problems determine how the loop should allocate oracle calls, reuse past feedback, shape rewards, encourage exploration, and preserve validity.
%We study this question through online fine-tuning for discrete-diffusion molecular optimization. Starting from a pretrained generator, the optimizer repeatedly samples molecules, decides which candidates to evaluate, receives oracle scores, and updates the generator from the scored molecules. This loop creates a growing history of feedback, but it leaves two central design problems open. The first is exploration: the optimizer must both query informative molecules to learn the unknown reward landscape and generate candidates beyond regions already favored by the diffusion prior. The second is objective mismatch: the oracle rewards only task-relevant molecules, while likelihood-based fine-tuning can spend capacity modeling low-value or invalid samples. These problems determine how the loop should allocate oracle calls, reuse past feedback, shape rewards, encourage exploration beyond the prior's dominant modes, and preserve validity.
We therefore perform an integrated study of online training loops for  molecular optimization using discrete diffusion. Prior work has studied the main ingredients of this loop, including inference-time search with a fixed generator \citep{didi2026scaling, wangFineTuningDiscreteDiffusion2024, singhalGeneralFrameworkInferencetime2025, liDerivativeFreeGuidanceContinuous2024, nisonoffUnlockingGuidanceDiscrete2025, gruverProteinDesignGuided2023,  xiongProteinGuideOntheflyProperty2026, luProVADAGenerationSubcellular2025, ueharaRewardGuidedIterativeRefinement2025, liDynamicSearchInferenceTime2025, jianGeneralBindingAffinity2026, yangHitLeadDiscovery, tangTR2D2TreeSearch2025, yangSteeringGenerativeModels2025}, online fine-tuning from scored samples \citep{rector-brooksSteeringMaskedDiscrete2024, suIterativeDistillationRewardGuided2025, ueharaRewardGuidedIterativeRefinement2025, yangSteeringGenerativeModels2025, stoccoGuidingGenerativeProtein2025, stoccoGuidingGenerativeModels2025, liCAGenMolConditionAwareDiffusion2026, olivecronaMolecularNovoDesign2017, blalockFunctionalAlignmentProtein2025}, and reward-guided diffusion objectives that turn oracle feedback into a training loss \citep{wangFineTuningDiscreteDiffusion2024, ueharaRewardGuidedIterativeRefinement2025, ueharaFeedbackEfficientOnline2024, santiFlowDensityControl, tangTR2D2TreeSearch2025, ueharaInferenceTimeAlignmentDiffusion2025}.  We ask how representative design choices around the objective interact: whether choices for acquisition, feedback reuse, reward shaping, prior debiasing, and validity control provide distinct value, reinforce one another, or become redundant once combined \needcite.

%These pressures motivate complementary feedback mechanisms. Reward-model uncertainty motivates acquisition rules that use uncertainty to decide which candidates receive expensive oracle evaluations. Bias in the generative prior motivates prior debiasing, which encourages exploration beyond regions already favored by the current generator. Best-design discovery motivates reward shaping before fine-tuning, so high-reward candidates exert more influence on the update. Feedback reuse motivates replay, so useful evaluated molecules are not discarded after one gradient step. Validity constraints motivate penalties that keep optimization from exploiting invalid strings. Rather than proposing a new core fine-tuning objective, we ask how these mechanisms interact around existing reward-guided diffusion objectives.

%We study design choices across the major stages of the loop: selecting candidates for oracle evaluation, reusing evaluated samples, shaping oracle scores before fine-tuning, encouraging the generator to move beyond regions already favored by its pretrained prior, and preserving validity under aggressive optimization. This framing lets us test whether mechanisms aimed at different pressures---reward-model uncertainty, generative-prior bias, best-design discovery, feedback retention, and validity---compose into a stronger online loop or merely duplicate one another.


%Online fine-tuning still leaves several choices open. After each batch is scored, the method must choose which new candidates are worth evaluating, which past candidates should remain in the training set, how much weight to put on the highest-scoring molecules, and how to keep the generator from producing invalid molecules. Emphasizing only the best molecules may improve top scores but reduce exploration. Pushing the model away from familiar regions may uncover new candidates but damage validity. Reusing past high-scoring molecules can stabilize learning, but may slow adaptation to new regions. This leaves an empirical question: which parts of the online loop matter most, do their benefits add or saturate, and how do they behave under shared oracle-call and runtime accounting?

%We turn these problems into a controlled study of the online feedback loop. To address statistical exploration, we use Thompson sampling (CITE) over a reward model so oracle calls target molecules that are both promising and uncertain. To address computational exploration, we add a model-debiasing term that discourages the generator from staying only in high probability regions of its current distribution. To align training with the optimization goal, we use a tail-focused objective that concentrates updates on high-scoring molecules. Because this pressure can push the generator toward invalid structures, we add an invalid-output penalty during fine-tuning. Finally, because oracle calls are costly, we maintain a replay buffer of evaluated molecules so useful feedback can influence multiple generator updates. We compare against offline fine-tuning to test whether feedback must be incorporated online, and against search-based variants to test whether additional inference-time exploration still helps once the generator itself is being updated.
Our experiments show that these mechanisms are complementary rather than redundant. In particular, likelihood-based debiasing improves exploration even when acquisition already accounts for reward-model uncertainty, suggesting that uncertainty in the reward model and bias in the generative prior are distinct bottlenecks. Reward shaping, replay, and validity control affect different parts of the training signal and failure profile rather than substituting for a stronger fine-tuning objective. Across 6 small-molecule binding-affinity and 3 protein fitness optimization tasks, the combined online training loop [RESULT] offline fine-tuning and search-augmented variants under matched oracle-call budgets, with [RESULT] under [GPU-hour / wall-clock] accounting.

%We study this question in discrete-diffusion molecular optimization by turning the online loop into a set of testable design choices. Starting from a pretrained generator, we repeatedly sample molecules, choose which ones receive oracle calls, store evaluated molecules for future updates, adjust their rewards, and fine-tune the generator. Each component targets one part of this loop. An exploration bonus on the reward model helps allocate oracle calls. A replay buffer retains useful feedback across updates. A tail-focused reward estimate makes learning less sensitive to low-value samples. A negative likelihood regularization term pushes the generator away from over-occupied regions, while an invalid-output penalty keeps fine-tuning near the valid molecular manifold. Offline fine-tuning and search-based variants serve as diagnostic baselines, testing whether feedback must be incorporated online and whether additional inference-time exploration still helps once the generator is being updated.

%Our experiments show that the main design choices are complementary rather than redundant: the strongest results require both forms of exploration, reward shaping, replay, and validity control. Each choice improves a different part of the online loop. Statistical exploration improves where oracle calls are spent, computational exploration expands the regions the generator reaches, reward shaping determines how scores become gradients, replay preserves useful feedback across updates, and validity control prevents reward optimization from producing invalid molecules. Ablating these choices reveals distinct failure modes rather than a single dominant component. Across [TASKS/TARGETS], the combined recipe [RESULT] offline fine-tuning and search-based variants under matched oracle-call budgets, with [RESULT] under [GPU-hour / wall-clock] accounting.

Our contributions are:
\begin{itemize}
    \item We formulate test-time molecular optimization as an online feedback-allocation problem, where a fixed budget must be divided across candidate evaluation, reward modeling, replay, reward shaping, and generator updates.

    \item We provide a controlled study of online fine-tuning for discrete-diffusion molecular optimization under oracle-call and runtime accounting, with offline fine-tuning and search-based variants included as diagnostic baselines.

    \item We show that several simple additions to the online loop provide complementary benefits: Thompson sampling over the reward model improves oracle-call allocation, replay stabilizes learning, CVaR reward shaping focuses updates on high-reward candidates, model debiasing encourages exploration beyond dominant modes, and invalid-output penalties preserve molecular validity.

    \item We combine these components into a single online active fine-tuning recipe that [RESULT] across [TASKS/TARGETS] under matched oracle-call budgets, with [RESULT] under [GPU-hour / wall-clock] accounting.
\end{itemize}